% Figure: a canal mitre lock gate seen in plan. Two leaves meet at a mitre
% pointing upstream; water pressure on the upstream face drives each leaf
% against its neighbour at the mitre and against the quoin (heel) post in the
% wall. Panel (a) plan of the pair; panel (b) FBD of one leaf as a three-force
% body. Numbers match the lock-gate vignette.
\begin{figure}[htbp]
\centering
\begin{tikzpicture}[
  force/.style={draw=engblue, -{Stealth}, very thick},
  rxn/.style={draw=engred, -{Stealth}, very thick},
  water/.style={draw=engblue!45, -{Stealth}, thick},
  scale=1.0,
]
% ----- panel (a): plan of the mitred pair -----
\begin{scope}[shift={(0,0)}]
% chamber walls
\draw[draw=enggray, line width=2pt] (-0.2,2.2) -- (-0.2,0.6);
\draw[draw=enggray, line width=2pt] (5.2,2.2) -- (5.2,0.6);
\node[font=\scriptsize, enggray, anchor=east] at (-0.25,1.9) {wall};
% quoin (heel) posts in the wall corners
\fill[enggray!35] (-0.2,0.6) rectangle (0.15,0.95);
\fill[enggray!35] (4.85,0.6) rectangle (5.2,0.95);
\node[font=\scriptsize, enggray, anchor=south east] at (0.15,0.95) {quoin};
% two leaves meeting at the mitre point M (pointing upstream, downward here)
\coordinate (Lh) at (0.0,0.8);   % left heel
\coordinate (Rh) at (5.0,0.8);   % right heel
\coordinate (M)  at (2.5,-0.15);  % mitre point, upstream
\draw[draw=engblack, line width=2.4pt] (Lh) -- (M);
\draw[draw=engblack, line width=2.4pt] (Rh) -- (M);
\fill[engblack] (M) circle (0.05);
\node[font=\scriptsize, anchor=north] at (M) {$M$};
\node[font=\scriptsize, anchor=south] at (Lh) {heel};
% upstream / downstream labels
\node[font=\scriptsize\itshape, engblue, anchor=north] at (2.5,-0.55) {upstream (high water)};
\node[font=\scriptsize\itshape, enggray, anchor=south] at (2.5,1.15) {downstream (low water)};
% water pressure arrows on the upstream face of the left leaf (pushing downstream)
\draw[water] (1.55,0.05) -- ++(0.0,0.55);
\draw[water] (1.05,0.28) -- ++(0.0,0.55);
\draw[water] (0.55,0.52) -- ++(0.0,0.55);
\node[font=\scriptsize, anchor=north east] at (0.55,0.05) {(a)};
\end{scope}

% ----- panel (b): FBD of one leaf as a three-force body -----
\begin{scope}[shift={(8.4,0.4)}]
\coordinate (H) at (0,0);      % heel (quoin) reaction point
\coordinate (Mp) at (2.5,-0.9); % mitre contact point
\draw[draw=engblack, line width=2.4pt] (H) -- (Mp);
\fill[engblack] (H) circle (0.05);
\fill[engblack] (Mp) circle (0.05);
\node[font=\scriptsize, anchor=south east] at (H) {heel $H$};
\node[font=\scriptsize, anchor=north west] at (Mp) {mitre $M$};
% resultant water thrust P, acting at the leaf mid-length, perpendicular to leaf
\coordinate (C) at (1.25,-0.45);
\draw[force] (C) -- ++(0.62,0.83) node[anchor=south west, font=\small, engblue] {$P$};
\node[font=\scriptsize, engblue, anchor=north east] at (C) {mid-span};
% heel reaction at H (component form): horizontal H_x and along-gate H_y
\draw[rxn] (H) -- ++(-1.0,0.35) node[anchor=east, font=\small, engred] {$R_H$};
% mitre reaction: normal to the contact face at M
\draw[rxn] (Mp) -- ++(1.05,0.30) node[anchor=west, font=\small, engred] {$R_M$};
\node[font=\scriptsize, anchor=north] at (1.25,-1.45) {(b)};
\end{scope}
\end{tikzpicture}
\caption[A mitre lock gate as a three-force body]{A canal lock's mitre gate,
plan view. Panel (a): two leaves meet at the mitre point $M$, which points
upstream so that the higher water presses the leaves together rather than
apart. Each leaf pivots about its heel (quoin) post set into the wall. The
upstream water pressure grows with depth, and its resultant $P$ acts
perpendicular to the leaf at the centroid of the triangular pressure
distribution. Panel (b): the free-body diagram of one leaf. Three forces act:
the water thrust $P$ at mid-span, the heel reaction $R_H$ where the leaf bears
in its quoin, and the mitre reaction $R_M$ where one leaf presses on the
other. With only three non-parallel forces, the leaf is a three-force body and
its lines of action are concurrent: extending the lines of $P$ and $R_M$ to
their crossing fixes the direction of $R_H$ before any component equation is
written. The gate is closed and held by water alone, which is why a mitre gate
needs no latch against the head it faces and why it must be opened against
almost no load once the levels are equalised.}
\label{fig:vol03ch01:lock-gate}
\end{figure}
